\chapter{Results}
\label{chap:findings}

This chapter reports how far an algorithm could take over the day's planning work, and where the coordinator's Control sets the limit. The interviews in \Cref{sec:interview-themes} say where coordinators welcome automation and where they will not give it up. The benchmark in \Cref{sec:algorithm-benchmark} reports how the available solvers perform on synthetic instances. The system description in \Cref{sec:engineering-results} shows what was built to operate within those bounds, and \Cref{sec:administrative-results} closes the chapter with the sprint record. RP was not user-tested with coordinators during the project period; implications for validation are taken up in Chapter 5.


\section{Scientific Results}
\label{sec:scientific-results}

\subsection{Interview Themes}
\label{sec:interview-findings}
\label{sec:interview-themes}

All seven interviewed companies plan the day's assignments manually. Their tools differ, from phone and memory to Timpex, Opter, and internal systems, but none of them generates the driver-vehicle match. The planning decision remains with the coordinator.

The most surprising theme was how little coordinators can see about resource utilisation while they plan. Overtime is handled when it surfaces, not tracked toward limits. Idle time between assignments is not visible in any of the tools the seven companies use. Load across drivers is balanced from the coordinator's memory of who has done what this week, not from a structured view.

Owners and Admmit asked for better utilisation visibility; coordinators asked for speed and override authority. The interviews therefore show a split between the people asking for optimisation and the people who must work with the suggested plans each day.

Two pain points were named across the seven calls: dependence on tacit knowledge, and the absence of a single capacity overview. Every coordinator carried operational knowledge no current system captures, and treated this as the reason a fully automated plan would be unsafe. No interviewed company has a single view that combines driver availability, vehicle status, and assignment load; coordinators reconstruct that picture each morning from several sources and from memory.

Sick-leave handling varies materially across the seven companies. Some coordinators describe it as a small task: Norlog Lakselv handles a sick-day in ``to-tre tastetrykk'' (two or three keystrokes), and Nordic Crane reports finishing a replanning event in about twenty-seven seconds. Others describe sick-leave as the dominant problem of their day. Norlog Tana put it directly: ``jeg sitter jo egentlig med den problemløsningen hver eneste dag'' (I sit with that problem-solving every single day). The difference depends on how many drivers are available, how flexible the routes are, and whether the routes need special competence such as scanning at Norlog Tana. Ottem describes a structurally similar problem from the order side: cancellations come in unpredictably and force the same kind of reshuffling.

Attitudes toward algorithm-assisted planning were split, not aligned. Norlog Lakselv, Norlog Mo i Rana, and Norlog Tana welcomed the idea. Harlem Solutions, Ottem, and Bergen Bulk Transport were sceptical, mainly because transport is unpredictable and they wanted to keep human judgement in the loop. Nordic Crane was partly positive, asking specifically for automatic distribution eighteen to twenty-four hours before an assignment runs, rather than full real-time automation. Even the sceptical interviewees accepted the idea of an algorithmic suggestion. Ottem framed the shared expectation most directly: ``det må være noen som sitter og tenker bak meg da'' (``there must be someone sitting and thinking behind me''). Across the seven calls, the agreement was that the system should propose a plan the coordinator can correct, not replace the coordinator.

When asked what they consider in matching a driver to an assignment, coordinators name the same criteria, in a consistent frequency order: availability and current position, experience (especially on demanding routes), driving and rest-time status (legally required), competency and licence, route familiarity (especially in northern Norway), equipment match, and overtime status toward annual limits. The criteria were stable across companies, but their weighting was not. Some prioritised workload balance, others vehicle-driver continuity or route familiarity.

\Cref{tab:themes} summarises the themes the analysis named, the frequency at which each surfaced across the seven calls, and the direction each theme pointed.

\begin{table}[h]
\centering
\caption{Themes named in the interview analysis. Frequency reports how widely each theme surfaced across the seven calls; direction reports whether the theme was unanimous or split.}
\label{tab:themes}
\small
\begin{tabular}{p{0.30\textwidth}p{0.28\textwidth}p{0.34\textwidth}}
\toprule
\textbf{Theme} & \textbf{Frequency} & \textbf{Direction} \\
\midrule
Resource-utilisation visibility gap & All 7 companies & Universal: overtime reactive, idle time invisible, load balance from memory \\
Operator-vs-owner asymmetry & Articulated by owners and Admmit, not by coordinators & Mirrors Bainbridge's ironies of automation \\
Tacit-knowledge dependency & All 7 coordinators & Universal: knowledge no current system captures, treated as the reason full automation is unsafe \\
Lack of capacity overview & All 7 companies & Universal: no single view combining driver availability, vehicle status, assignment load \\
Sick-leave handling & All 7 (severity varies) & Mixed: small task (Norlog Lakselv ``to-tre tastetrykk'') to dominant (Norlog Tana ``hver eneste dag'') \\
Attitudes to automation & 3 positive, 3 sceptical, 1 partial & Split, but consensus on ``suggest, don't decide'' \\
\bottomrule
\end{tabular}
\end{table}

\subsection{Algorithm Benchmark}
\label{sec:algorithm}
\label{sec:algorithm-benchmark}

The benchmark shows how the solver engines behave on the same synthetic planning instances under the protocol described in \Cref{sec:eval-benchmark}. The result is a controlled comparison of solver behaviour across small, medium, and large resource-constrained datasets.

Each benchmark assignment needs both a driver and a vehicle. The pair must satisfy the operational hard constraints for availability, competence, vehicle suitability, and overlapping time windows. Assignments for which no feasible pair exists remain unassigned with a recorded reason.

The seed-42 benchmark shows that no single engine dominates all instance sizes. Greedy is the fastest baseline and returns a feasible solution on all three datasets, but schedules fewer assignments than the stronger search engines where those complete. OR-Tools CP-SAT schedules the most assignments on the small and medium datasets and proves optimality in its own model, but times out on the large dataset under the 120-second cap. Timefold underperforms on the small dataset, where CP-SAT reaches the proven optimum almost immediately, but improves over greedy on the medium and large datasets and is the only non-greedy engine to complete the large instance within the cap.

\begin{table}[h]
\centering
\caption{Multi-engine benchmark, seed-42 datasets, 120-second per-engine cap. OR-Tools CP-SAT proves optimality on small and medium but times out on large; Timefold is the only engine that completes the large instance with a coverage improvement over greedy.}
\label{tab:benchmark-results}
\small
\begin{tabular}{llll}
\toprule
\textbf{Dataset} & \textbf{Greedy} & \textbf{OR-Tools CP-SAT} & \textbf{Timefold} \\
\midrule
Small  (20 assignments)  & 10/20 (50.0\,\%)   & 11/20 (55.0\,\%, optimal)   & 9/20 (45.0\,\%) \\
\midrule
Medium (100 assignments) & 47/100 (47.0\,\%)  & 67/100 (67.0\,\%, optimal)  & 60/100 (60.0\,\%) \\
\midrule
Large  (500 assignments) & 189/500 (37.8\,\%) & timeout ($>$120\,s)         & 222/500 (44.4\,\%) \\
\bottomrule
\end{tabular}
\end{table}

All successful runs have zero hard-tier violations, which makes the coverage figures comparable across engines. The canonical scorer still reports legal-tier violations in every successful solution. These violations concern working-time and rest-period rules, such as minimum rest, weekly hours, and consecutive working days. They grow with dataset size because tight resource pools force longer employee schedules. In RP, those issues are surfaced for coordinator review rather than treated as hard infeasibility by the benchmark engines.

The normal solve flow uses these differences without asking the coordinator to choose an engine. The coordinator chooses \texttt{rask}, \texttt{middels}, or \texttt{grundig}. RP then decides which engine to run from the allowed pool for that category, using earlier runs when enough history exists and trying the available engines when it does not. The older explicit \texttt{engineId} mode remains available for benchmark and advanced use, but it is not the ordinary coordinator flow.

\section{Engineering Results}
\label{sec:engineering-results}

\subsection{System Overview}
\label{sec:system-description}
\label{sec:requirements}

RP is the implemented artefact built around the planning step identified in \Cref{sec:interview-themes}. It is a web application where the coordinator works with company-scoped planning data and asks solver engines for proposed plans. At result level, RP consists of one web application, one PostgreSQL database, and three solver engines. Each company's data is separated by \texttt{companyId}, so one deployment can serve several companies without mixing their assignments, employees, vehicles, planning preferences, or optimisation history. \Cref{fig:architecture} shows this structure.

\begin{figure}[h]
  \centering
  \resizebox{\textwidth}{!}{%
  \begin{tikzpicture}[
    node distance=0.65cm,
    box/.style={rectangle, draw, rounded corners, align=center, minimum height=0.85cm, font=\small, fill=blue!4},
    wide/.style={box, minimum width=11cm},
    solver/.style={box, minimum width=3.4cm, fill=orange!7},
    db/.style={box, minimum width=4cm, fill=green!6},
    arrow/.style={-Stealth, thick}
  ]
    \node[wide] (browser) {Browser client};
    \node[wide, below=of browser] (server) {Web application\\\footnotesize coordinator screens and business logic};
    \node[db, below=1.2cm of server, xshift=-3.5cm] (db) {PostgreSQL\\\footnotesize multi-tenant via \texttt{companyId} on every row};
    \node[box, below=1.2cm of server, xshift=3cm, minimum width=4.5cm, fill=orange!7] (registry) {Solver registry\\\footnotesize shared problem format};
    \node[solver, below=of registry, xshift=-3.6cm] (greedy) {Greedy heuristic\\\footnotesize Python subprocess};
    \node[solver, below=of registry] (ortools) {OR-Tools CP-SAT\\\footnotesize Python subprocess};
    \node[solver, below=of registry, xshift=3.6cm] (timefold) {Timefold metaheuristic\\\footnotesize JVM subprocess};

    \draw[arrow] (browser) -- (server);
    \draw[arrow] (server) -- (db);
    \draw[arrow] (server) -- (registry);
    \draw[arrow] (registry) -- (greedy);
    \draw[arrow] (registry) -- (ortools);
    \draw[arrow] (registry) -- (timefold);
  \end{tikzpicture}%
  }
  \caption{System architecture of RP. One web application serves the coordinator interface and business logic, PostgreSQL stores company-scoped data, and three solver engines return candidate plans through a shared problem format.}
  \label{fig:architecture}
\end{figure}

The web application holds the coordinator screens and server logic. The database stores company-scoped planning data and settings. The solver engines run outside the web application and return candidate plans to the same planning interface.

RP was specified against 58 requirements: 42 functional requirements and 16 non-functional requirements, prioritised under MoSCoW. The full registers are placed in \Cref{app:requirements-functional,app:requirements-nonfunctional}, and the traceability check is described in \Cref{sec:eval-traceability}. The workflow and algorithm results below show what those requirements became in the implemented artefact.

\subsection{Coordinator Workflow}
\label{sec:coordinator-workflow}

The coordinator's main screen is the planning timeline. It shows assignments across the selected planning period, with employees and vehicles available as planning resources. From that view, the coordinator can read the day, adjust assignments, and approve the plan.

When the coordinator drags an assignment to another driver or vehicle, RP saves the change and continues even if the move creates a conflict. The conflict is shown as a warning on the timeline and in the deviation viewer. The system records the problem without blocking the coordinator's decision.

Plan generation starts from a time-quality category: \texttt{rask}, \texttt{middels}, or \texttt{grundig}. After the run, RP shows the proposed plan with its score, violations, chosen engine, and whether the engine choice came from previous run history or from trying the available engines. The coordinator can approve the proposal, change it, or run the solver again.

\subsection{The Algorithm}
\label{sec:algorithm-implementation}

The algorithm turns the coordinator's planning problem into a set of driver-vehicle choices. The input is a planning period with assignments, employees, and vehicles. The output is a set of assignments with the employees and vehicles allocated to each one. Each possible allocation is either allowed or not allowed, and the solver searches for the set of allowed allocations that gives the best plan.

All three solver engines receive the same planning problem and return the same kind of solution. The difference is search strategy: greedy builds one plan quickly, CP-SAT searches systematically within a time limit, and Timefold improves a candidate plan over time. The benchmark result in \Cref{sec:algorithm-benchmark} reports how those differences behaved on the seed-42 datasets.

Normal plan generation asks the coordinator to choose a category, not an engine. The category defines an engine pool and a time cap, and the auto-router decides which engine to run from that pool. When there is not enough history, it tries the available engines. When there is enough history, it runs the engine that has performed best for that company, size bucket, and category. An explicit \texttt{engineId = "auto"} mode remains available as an engine-portfolio fallback, but it is not the normal coordinator flow.

A driver without the required competence cannot be assigned just because the plan looks good elsewhere. RP therefore treats competence, availability, vehicle status, vehicle type, and double-booking as hard constraints.

\begin{table}[h]
\centering
\caption{Hard constraints. A plan that violates any of these is not feasible.}
\label{tab:hard-constraints}
\small
\begin{tabularx}{\textwidth}{@{}
  p{0.14\textwidth}
  X
@{}}
\toprule
\textbf{ID} & \textbf{Constraint} \\
\midrule
HC-01 & Employee competencies match the assignment requirement \\
HC-02 & Employee is available for the assignment time \\
HC-03 & Vehicle is available and active \\
HC-04 & Vehicle type matches the assignment requirement \\
HC-05 & No employee is double-booked across overlapping assignments \\
HC-06 & No vehicle is double-booked across overlapping assignments \\
\bottomrule
\end{tabularx}
\end{table}

After feasibility is secured, RP can still prefer one valid plan over another. The soft constraints express those preferences, with per-company weights for continuity, balance, priority, transitions, and employee preferences.

\begin{table}[h]
\centering
\caption{Soft constraints. Each carries a per-company weight.}
\label{tab:soft-constraints}
\small
\begin{tabularx}{\textwidth}{@{}
  p{0.14\textwidth}
  X
  p{0.18\textwidth}
@{}}
\toprule
\textbf{ID} & \textbf{Constraint} & \textbf{Weight} \\
\midrule
SC-01 & Prefer assigning an employee to their designated vehicle & Configurable \\
SC-02 & Balance workload across employees & Configurable \\
SC-03 & Minimise travel or transition cost between consecutive assignments & Configurable \\
SC-04 & Prioritise high-priority assignments & Configurable \\
SC-05 & Respect stated employee preferences & Configurable \\
\bottomrule
\end{tabularx}
\end{table}

A plan that breaks a strict rule should not beat a valid plan because it performs better on softer goals. RP therefore scores plans in tiers: hard violations first, legal working-time rules next, then business-critical quality, and finally preference quality. Coverage is rewarded inside the business-critical tier, after the stricter rules have been handled.

\section{Administrative Results}
\label{sec:administrative-results}

\subsection{Project Timeline}
\label{sec:sprint-summary}

The project ran from 12 January to 9 April 2026 as seven calendar sprints, with thesis writing continuing toward submission in May. The sprint log records project administration and delivery rhythm; the six learning iterations in \Cref{sec:iterations} record what the artefact learned technically and empirically. The two should therefore be read as different views of the same work: sprints describe time, iterations describe learning.

\Cref{tab:sprint-timeline} summarises the sprint record. The project started with scope, stack choice, and project setup, then moved through schema, authentication, deviation handling, solver engines, interviews, integration, benchmarking, and documentation.

\begin{table}[h]
\centering
\caption{Sprint timeline for the RP project.}
\label{tab:sprint-timeline}
\small
\begin{tabularx}{\textwidth}{@{}
  p{0.16\textwidth}
  p{0.22\textwidth}
  X
@{}}
\toprule
\textbf{Sprint} & \textbf{Period} & \textbf{Main work} \\
\midrule
Sprint 0 & 12--16 Jan & Kickoff, scope, stack choice, project setup \\
Sprint 1 & 19--30 Jan & Scaffold, schema, authentication, database integration \\
Sprint 2 & 2--13 Feb & Deviation handling, sick-leave flow, interface polish \\
Sprint 3 & 16--27 Feb & Greedy, OR-Tools CP-SAT, and Timefold engines \\
Sprint 4 & 2--13 Mar & Seven interviews, fit/gap work, user-research features \\
Sprint 5 & 16--27 Mar & Integration, end-to-end tests, thesis chapter drafting \\
Sprint 6 & 30 Mar--9 Apr & Benchmarks, hardening, documentation \\
\bottomrule
\end{tabularx}
\end{table}

These results provide the empirical, technical, and administrative basis for the discussion of Efficiency, Control, and Adaptability in Chapter 5.
